\chapter{Surfaces}\label{ch:b2:surfaces}

After curves, surfaces: two-parameter objects in $\R^3$. The
differential calculus of \cref{ch:b2:diffcalc} provides everything
we need --- partial derivatives give tangent vectors, the cross
product gives the normal, determinants give areas. We define
regular parametrized surfaces, their tangent planes, and the
\emph{first fundamental form}, which encodes all length and area
measurements on the surface. Surfaces also arise as level sets
$f(x, y, z) = c$; the gradient then directs the normal.

\section{Parametrized surfaces}

\begin{definition}[Regular parametrized surface]\label{def:b2:surfaces:param}
Let $U \subseteq \R^2$ be open. A \emph{parametrized surface} of
class $\mathcal{C}^k$ ($k \geq 1$) is a map $\sigma \colon U \to
\R^3$, $(u, v) \mapsto \sigma(u, v)$, of class $\mathcal{C}^k$. A
point is \emph{regular} if the partial derivative vectors
\[
\sigma_u(u, v) = \frac{\partial\sigma}{\partial u}(u,v),
\qquad
\sigma_v(u, v) = \frac{\partial\sigma}{\partial v}(u,v)
\]
are linearly independent, i.e.\ $\sigma_u \wedge \sigma_v \neq 0$;
the surface is \emph{regular} if every point is.
\end{definition}

\begin{example}[The three standard descriptions]\label{ex:b2:surfaces:standard}
\leavevmode
\begin{enumerate}
\item \emph{Graph:} $\sigma(u, v) = (u,\ v,\ f(u, v))$ for $f \in
\mathcal{C}^1(U)$. Always regular: $\sigma_u = (1, 0, f_u)$ and
$\sigma_v = (0, 1, f_v)$ are independent.
\item \emph{Sphere (spherical coordinates):} for the sphere of
radius $R$,
\[
\sigma(\theta, \varphi)
= (R\cos\theta\cos\varphi,\ R\sin\theta\cos\varphi,\
R\sin\varphi),
\qquad
(\theta, \varphi) \in \R \times
\bigl(-\tfrac\pi2, \tfrac\pi2\bigr),
\]
with $\theta$ the longitude and $\varphi$ the latitude. One checks
$\norm{\sigma_\theta \wedge \sigma_\varphi} = R^2\cos\varphi > 0$:
regular away from the poles (which this chart omits).
\item \emph{Level set:} $S = \{(x,y,z) : f(x, y, z) = c\}$ where
$f$ is $\mathcal{C}^1$ and $\nabla f \neq 0$ on $S$. Near each
point, one coordinate can be expressed as a function of the other
two by the implicit function theorem (\cref{ch:b2:diffcalc}), so
$S$ is locally a graph.
\end{enumerate}
\end{example}

\section{Tangent plane and normal}

\begin{definition}[Tangent plane]\label{def:b2:surfaces:tangent}
Let $\sigma$ be regular at $(u_0, v_0)$, $M_0 = \sigma(u_0, v_0)$.
The \emph{tangent plane} $T_{M_0}S$ is the plane through $M_0$
directed by $\operatorname{Vect}(\sigma_u, \sigma_v)$ (partials at
$(u_0, v_0)$). The \emph{unit normal} is
\[
n(u_0, v_0)
= \frac{\sigma_u \wedge \sigma_v}
       {\norm{\sigma_u \wedge \sigma_v}} .
\]
\end{definition}

\begin{proposition}[Tangent vectors are velocity
vectors]\label{prop:b2:surfaces:velocity}
The direction of $T_{M_0}S$ is exactly the set of vectors
$\gamma'(0)$, where $\gamma = \sigma \circ c$ ranges over the
$\mathcal{C}^1$ curves drawn on the surface through $M_0$ (i.e.\
$c \colon (-\varepsilon, \varepsilon) \to U$ is $\mathcal{C}^1$
with $c(0) = (u_0, v_0)$).
\end{proposition}

\begin{proof}
If $c(t) = (u(t), v(t))$, the chain rule
(\cref{ch:b2:diffcalc}) gives
\[
\gamma'(0) = u'(0)\,\sigma_u + v'(0)\,\sigma_v
\in \operatorname{Vect}(\sigma_u, \sigma_v).
\]
Conversely, the vector $a\sigma_u + b\sigma_v$ is attained by the
curve $c(t) = (u_0 + at,\ v_0 + bt)$, which stays in the open set
$U$ for $\abs t$ small.
\end{proof}

\begin{proposition}[Normal of a level surface]\label{prop:b2:surfaces:gradient}
Let $S = \{f = c\}$ with $f$ of class $\mathcal{C}^1$ and $\nabla
f(M_0) \neq 0$. Then the tangent plane of $S$ at $M_0$ is the
plane through $M_0$ orthogonal to $\nabla f(M_0)$:
\[
T_{M_0}S :\quad
\langle \nabla f(M_0),\ M - M_0\rangle = 0 .
\]
\end{proposition}

\begin{proof}
For any curve $\gamma$ drawn on $S$ through $M_0$, $f(\gamma(t)) =
c$ identically, so the chain rule gives $\langle
\nabla f(M_0), \gamma'(0)\rangle = 0$: all velocity vectors are
orthogonal to the gradient, so the tangent direction is contained
in the plane $\nabla f(M_0)^\perp$. Both are $2$-dimensional
subspaces --- the tangent direction because $S$ is locally a
regular graph (\cref{ex:b2:surfaces:standard}), the orthogonal
complement because $\nabla f(M_0) \neq 0$ --- hence they are
equal.
\end{proof}

\begin{example}\label{ex:b2:surfaces:spheretangent}
For the sphere $x^2 + y^2 + z^2 = R^2$: $\nabla f = 2(x, y, z)$,
so the tangent plane at $M_0$ is orthogonal to the radius
$\vect{OM_0}$ --- the classical fact that radius and tangent plane
are perpendicular, with equation $\langle M_0, M\rangle = R^2$.
\end{example}

\begin{example}[Graph tangent plane]\label{ex:b2:surfaces:graphtangent}
For $z = f(x, y)$ at $(x_0, y_0)$: applying
\cref{prop:b2:surfaces:gradient} to $F(x,y,z) = f(x,y) - z$,
\[
z = f(x_0, y_0)
+ f_x(x_0, y_0)(x - x_0)
+ f_y(x_0, y_0)(y - y_0) ,
\]
the affine part of the first-order Taylor expansion --- the
tangent plane is the graph of the differential, as it must be.
\end{example}

\section{The first fundamental form}

\begin{definition}[First fundamental form]\label{def:b2:surfaces:fff}
Let $\sigma \colon U \to \R^3$ be a regular $\mathcal{C}^1$
surface. Its \emph{first fundamental form} at $(u,v)$ is the
positive definite quadratic form on $\R^2$
\[
I(h, k) = \norm{h\,\sigma_u + k\,\sigma_v}^2
= E\,h^2 + 2F\,hk + G\,k^2,
\]
where
\[
E = \norm{\sigma_u}^2,
\qquad
F = \langle \sigma_u, \sigma_v\rangle,
\qquad
G = \norm{\sigma_v}^2 .
\]
\end{definition}

\begin{remark}
$I$ is the restriction of the ambient Euclidean scalar product to
the tangent plane, read in the basis $(\sigma_u, \sigma_v)$: it is
positive definite precisely because $\sigma_u, \sigma_v$ are
independent (\cref{ch:b2:quadratic}). Every metric quantity on the
surface --- lengths of drawn curves, angles between them, areas ---
is computed from $E, F, G$ alone. Two surfaces with the same $E,
F, G$ in suitable parameters are \emph{isometric} even if they sit
differently in space: this is the starting point of intrinsic
geometry.
\end{remark}

\begin{proposition}[Length of a curve drawn on a
surface]\label{prop:b2:surfaces:curvelength}
If $\gamma(t) = \sigma(u(t), v(t))$, $t \in [a, b]$, is
$\mathcal{C}^1$, then
\[
L(\gamma) = \int_a^b
\sqrt{E\,u'^2 + 2F\,u'v' + G\,v'^2}\;\dd t ,
\]
with $E, F, G$ evaluated at $(u(t), v(t))$.
\end{proposition}

\begin{proof}
$\gamma' = u'\sigma_u + v'\sigma_v$ by the chain rule, so
$\norm{\gamma'}^2 = I(u', v') = Eu'^2 + 2Fu'v' + Gv'^2$; integrate
$\norm{\gamma'}$ (\cref{def:b2:curves:length}).
\end{proof}

\begin{lemma}[Lagrange identity]\label{lem:b2:surfaces:lagrange}
For all $a, b \in \R^3$:
$\norm{a \wedge b}^2 = \norm a^2 \norm b^2 - \langle a, b\rangle^2$.
In particular
\[
\norm{\sigma_u \wedge \sigma_v} = \sqrt{EG - F^2}.
\]
\end{lemma}

\begin{proof}
Both sides are unchanged if we replace $b$ by its component $b_\perp
= b - \frac{\langle a, b\rangle}{\norm a^2}a$ orthogonal to $a$
(for $a \neq 0$; the case $a = 0$ is trivial): the left side
because $a \wedge a = 0$, the right side by expanding $\norm
{b_\perp}^2 = \norm b^2 - \frac{\langle a,b\rangle^2}{\norm a^2}$
and $\langle a, b_\perp\rangle = 0$. So it suffices to prove the
identity for orthogonal $a, b$, where it reads $\norm{a \wedge
b} = \norm a \norm b$: true, since for orthogonal vectors the
cross product has norm $\norm a\norm b\,\abs{\sin\frac\pi2}$. The
displayed formula is the case $a = \sigma_u$, $b = \sigma_v$.
\end{proof}

\begin{definition}[Area]\label{def:b2:surfaces:area}
Let $\sigma \colon U \to \R^3$ be a regular injective
$\mathcal{C}^1$ surface and $K \subseteq U$ a compact domain on
which double integrals make sense (\cref{ch:b2:multint}). The
\emph{area} of the piece $\sigma(K)$ is
\[
\mathcal{A}
= \iint_K \norm{\sigma_u \wedge \sigma_v}\,\dd u\,\dd v
= \iint_K \sqrt{EG - F^2}\;\dd u\,\dd v .
\]
\end{definition}

\begin{remark}[Why this formula]
The rectangle $[u, u + \dd u] \times [v, v + \dd v]$ is mapped, to
first order, onto the parallelogram spanned by $\sigma_u\,\dd u$
and $\sigma_v\,\dd v$, whose area is $\norm{\sigma_u \wedge
\sigma_v}\,\dd u\,\dd v$: the definition integrates the local
area-distortion factor, exactly as arc length integrates the local
speed. Consistency with change of parameters is
\cref{exo:b2:surfaces:8}; consistency with the change-of-variables
formula for double integrals is discussed in
\cref{ch:b2:multint}.
\end{remark}

\begin{example}[Area of the sphere]\label{ex:b2:surfaces:spherearea}
For the spherical chart of \cref{ex:b2:surfaces:standard}:
\[
\sigma_\theta = R(-\sin\theta\cos\varphi,\
\cos\theta\cos\varphi,\ 0),
\qquad
\sigma_\varphi = R(-\cos\theta\sin\varphi,\
-\sin\theta\sin\varphi,\ \cos\varphi),
\]
so $E = R^2\cos^2\varphi$, $F = 0$, $G = R^2$ and $\sqrt{EG - F^2}
= R^2\cos\varphi$. Hence
\[
\mathcal{A}
= \int_{-\pi/2}^{\pi/2}\!\!\int_0^{2\pi}
R^2\cos\varphi\;\dd\theta\,\dd\varphi
= 2\pi R^2\,\bigl[\sin\varphi\bigr]_{-\pi/2}^{\pi/2}
= \boxed{4\pi R^2} .
\]
\end{example}

\begin{example}[Surface of revolution]\label{ex:b2:surfaces:revolution}
Rotate the curve $z \mapsto (r(z), 0, z)$, $r > 0$ of class
$\mathcal{C}^1$, around the $z$-axis:
\[
\sigma(\theta, z) = (r(z)\cos\theta,\ r(z)\sin\theta,\ z).
\]
Then $E = r(z)^2$, $F = 0$, $G = 1 + r'(z)^2$, so
\[
\mathcal{A} = \int_a^b\!\!\int_0^{2\pi}
r(z)\sqrt{1 + r'(z)^2}\;\dd\theta\,\dd z
= 2\pi\int_a^b r(z)\sqrt{1 + r'(z)^2}\,\dd z ,
\]
the classical formula (circumference $2\pi r$ times the slant
length element). For the cone $r(z) = kz$, $z \in [0, h]$: $\mathcal
A = 2\pi k\sqrt{1 + k^2}\,\frac{h^2}{2} = \pi \rho \ell$ with
$\rho = kh$ the base radius and $\ell = h\sqrt{1 + k^2}$ the slant
height --- the school formula, now derived rather than admitted.
\end{example}

\begin{figure}[ht]
\centering
\begin{tikzpicture}[scale=1.05]
  % A saddle surface z = x^2 - y^2 drawn as a wire mesh in oblique projection
  % projection: (x,y,z) -> (x + 0.45*y, z + 0.35*y)
  \begin{scope}
    % u-lines (x varies, y fixed)
    \foreach \yy in {-1,-0.5,0,0.5,1} {
      \draw[ocDarkBlue!70, thick, domain=-1:1, smooth, samples=25]
        plot ({\x + 0.45*\yy}, {\x*\x - \yy*\yy + 0.35*\yy});
    }
    % v-lines (y varies, x fixed)
    \foreach \xx in {-1,-0.5,0,0.5,1} {
      \draw[ocMint!80!black, thick, domain=-1:1, smooth, samples=25]
        plot ({\xx + 0.45*\x}, {\xx*\xx - \x*\x + 0.35*\x});
    }
    % tangent plane at origin: z = 0 -> parallelogram
    \draw[ocRed, dashed]
      (-0.8 - 0.36, -0.28) -- (0.8 - 0.36, -0.28)
      -- (0.8 + 0.36, 0.28) -- (-0.8 + 0.36, 0.28) -- cycle;
    \fill (0,0) circle (0.045) node[below right, font=\small] {$M_0$};
    % normal vector at origin
    \draw[->, thick] (0,0) -- (0,0.9)
      node[right, font=\small] {$n$};
  \end{scope}
\end{tikzpicture}
\caption{The saddle $z = x^2 - y^2$ near the origin, with its
coordinate curves ($u$-curves in blue, $v$-curves in green), the
tangent plane at $M_0 = (0,0,0)$ (dashed) and the unit normal $n$.
The surface crosses its tangent plane --- the two-dimensional
analogue of an inflection.}
\label{fig:b2:surfaces:saddle}
\end{figure}

\section{Exercises}

\begin{exercise}[$\star$]\label{exo:b2:surfaces:1}
Show that the tangent planes of the cone $z = \sqrt{x^2 + y^2}$
(minus its apex) all pass through the apex. \emph{(Parametrize by
$\sigma(\theta, r) = (r\cos\theta, r\sin\theta, r)$, $r > 0$.)}
\end{exercise}

\begin{exercise}[$\star$]\label{exo:b2:surfaces:2}
Find the tangent plane of the ellipsoid $\frac{x^2}{a^2} +
\frac{y^2}{b^2} + \frac{z^2}{c^2} = 1$ at a point $(x_0, y_0,
z_0)$ of the surface.
\end{exercise}

\begin{exercise}[$\star$]\label{exo:b2:surfaces:3}
Compute $E, F, G$ for the helicoid $\sigma(u, v) = (v\cos u,\
v\sin u,\ au)$, $a > 0$, and the area of the piece $0 \leq u \leq
2\pi$, $0 \leq v \leq 1$, as an integral (evaluate it using
$\int\sqrt{v^2 + a^2}\,\dd v = \frac12\bigl(v\sqrt{v^2+a^2} +
a^2\ln(v + \sqrt{v^2 + a^2})\bigr) + C$).
\end{exercise}

\begin{exercise}[$\star\star$]\label{exo:b2:surfaces:4}
(Torus) Parametrize the torus obtained by rotating the circle of
center $(R, 0, 0)$ and radius $r < R$ in the $xz$-plane around the
$z$-axis:
\[
\sigma(\theta, \psi) = \bigl((R + r\cos\psi)\cos\theta,\
(R + r\cos\psi)\sin\theta,\ r\sin\psi\bigr).
\]
Compute $E, F, G$, check regularity, and show that the area is
$4\pi^2 R r$ (Pappus: mean circumference $2\pi R$ times circle
length $2\pi r$).
\end{exercise}

\begin{exercise}[$\star\star$]\label{exo:b2:surfaces:5}
Show that the area of the graph of $f \in \mathcal{C}^1(K)$ is
$\iint_K \sqrt{1 + f_x^2 + f_y^2}\,\dd x\,\dd y$, and compute it
for the paraboloid piece $z = \frac12(x^2 + y^2)$ over the disk
$x^2 + y^2 \leq 1$ (polar coordinates,
\cref{ch:b2:multint}).
\end{exercise}

\begin{exercise}[$\star\star$]\label{exo:b2:surfaces:6}
A drawn curve $\gamma(t) = \sigma(u(t), v(t))$ on the sphere of
radius $R$ (spherical chart) has $u = \theta(t)$, $v =
\varphi(t)$. Write its length as an integral in $\theta, \varphi$
and prove that among curves joining two points of the same
meridian $\theta = \theta_0$, the meridian arc is the shortest.
\emph{(Bound the integrand below by $R\abs{\varphi'}$.)}
\end{exercise}

\begin{exercise}[$\star\star\star$]\label{exo:b2:surfaces:7}
(Normal lines of a sphere) Let $S$ be a regular level surface
$\{f = c\}$, connected, all of whose normal lines pass through a
fixed point $\Omega$. Show that $S$ is contained in a sphere
centered at $\Omega$. \emph{(Show that $\norm{M - \Omega}^2$ has
zero derivative along every curve drawn on $S$.)}
\end{exercise}

\begin{exercise}[$\star\star\star$]\label{exo:b2:surfaces:8}
(Area is geometric) Let $\Phi \colon U' \to U$ be a $\mathcal{C}^1$
diffeomorphism between open sets of $\R^2$ and $\tilde\sigma =
\sigma \circ \Phi$. Show that
\[
\norm{\tilde\sigma_{u'} \wedge \tilde\sigma_{v'}}
= \abs{\det J_\Phi}\,
\norm{(\sigma_u \wedge \sigma_v)\circ\Phi} ,
\]
and deduce, using the change-of-variables formula of
\cref{ch:b2:multint}, that the area of
\cref{def:b2:surfaces:area} does not depend on the chosen regular
parametrization.
\end{exercise}
